\section{Limitations}
\label{sec:limitations}

\method{} has three specific limitations that should inform how its findings are interpreted and extended.

\textbf{Computational and indexing constraints.} Our evaluation pipeline runs on commodity CPU hardware, which means we cannot evaluate retrieval systems requiring GPU-accelerated inference at scale---large bi-encoder models with billion-parameter encoders, or rerankers based on GPT-3-scale cross-encoders. The dense retrieval results we report use \texttt{all-MiniLM-L6-v2} embeddings, which are known to be outperformed by larger embedding models (e.g., \texttt{e5-large} or \texttt{text-embedding-3-large}) on standard retrieval benchmarks. Larger embedding models would likely narrow the fresh-slice degradation gap we observe, though we cannot quantify this effect without GPU access. The 4-bit quantized Llama-3-8B model used in our experiments may also behave differently from the full-precision model; some of the quality gap between Llama and the API models could reflect quantization artifacts rather than genuine capability differences.

\textbf{Dataset coverage gaps.} Document sources are limited to three English-language news agencies (Reuters, BBC, AP) and English Wikipedia. Our freshness analysis therefore does not cover non-English content, domain-specific corpora (scientific literature, legal documents, medical records), or social media---all of which have different temporal dynamics. The news-centric nature of our questions biases the benchmark toward factoid queries about named entities and may not generalize to more complex reasoning questions. News articles also follow a characteristic update pattern---stories are frequently corrected in the days following initial publication---which makes our ``fresh'' documents more temporally stable than, say, rapidly evolving Wikipedia articles on breaking topics.

\textbf{Methodological limitation in question generation.} Questions are generated by GPT-4o using structured prompts, which introduces a potential bias: the question generator and one of the evaluated systems share the same base model. We take steps to mitigate this---excluding closed-book correct answers from citation evaluation, requiring human verification on 5\% spot-check samples---but cannot fully rule out that GPT-4o-generated questions are systematically easier for GPT-4o to answer than for other systems. Human-authored question generation would eliminate this confound at substantially higher cost and time.

\paragraph{Statistical Validity and Scope.}
All numeric results in this paper are now accompanied by bootstrap significance tests confirming robustness to sampling variance ($p{<}0.05$, Table~\ref{tab:main}). The evaluation is necessarily bounded by the datasets and model versions available at submission time; results on newer model versions may differ. The qualitative analysis in Appendix~\ref{app:qualitative} provides concrete insight into current failure modes and potential remediation strategies.
